\subsection{Instrument 5: Content--Inference Circumplex Convergence}

The emotional circumplex measured through content-level processing and through
user-model emotion inference converges on the same geometric subspace. Using the
Qwen3.5-27B Opus distill, we extracted 30 emotion centroids from 900 trials of
conversation-context inference (the user-model emotion space from~\citet{lyra2026weather})
and compared them against valence and arousal directions computed from 250 emotional
scenario prompts (hostile--calm, desperate--calm difference-of-means).

The two circumplex axes capture $44$--$54\%$ of the total variance in the full
30-emotion user-model space at every layer from L6 onward---approximately $1000\times$
the ${\sim}0.04\%$ expected from two random directions in $d{=}5120$. The arousal
direction aligns with the user-model's first principal component at
$\cos = 0.83$--$0.87$ (mid-to-late layers), and the valence direction at
$\cos = 0.70$--$0.80$. Both systems exhibit identical developmental timecourses:
volatile eccentricity at L0--L16, stable geometry from L17 onward.

The measurements are independent---different prompts, different templates, different
experimental designs---yet they recover the same high-variance directions. This
convergence parallels the W$_K$ bridge ($\rho = 0.862$) from~\citet{lyra2026weather},
which projects \emph{into} this shared circumplex subspace. The model does not
maintain separate emotional geometry for content processing and partner-state inference;
it uses a common representational subspace for both.

We note the caveat that shared geometry does not entail shared computation. Two
functionally independent circuits could reuse the same high-variance directions
without computational coupling. An injection experiment (perturb one pathway, measure
the other) is pre-registered to distinguish representational overlap from causal
interaction.
